\documentclass[10pt, conference]{IEEEtran}
\IEEEoverridecommandlockouts
\usepackage{xcolor}        
\usepackage{graphicx}
\usepackage{amsmath, amssymb, amsfonts, amsthm}
\usepackage[caption=false, font=footnotesize]{subfig}
\usepackage{textcomp}
\usepackage{cite}
\usepackage{multirow}
\usepackage[ruled,vlined,linesnumbered]{algorithm2e}
\setlength{\algomargin}{1.5em}

%\allowdisplaybreaks

\DeclareMathOperator*{\argmax}{arg\,max}
\DeclareMathOperator*{\argmin}{arg\,min}
\DeclareMathOperator*{\maximize}{maximize}
\DeclareMathOperator*{\minimize}{minimize}
\DeclareMathOperator{\diag}{diag}
\DeclareMathOperator{\blkdiag}{blkdiag}
\DeclareMathOperator{\tr}{tr}
\DeclareMathOperator{\tril}{tril}
\DeclareMathOperator{\re}{Re}
\newtheorem{lemma}{Lemma}
\newtheorem{remark}{Remark}
\newcommand{\quotes}[1]{\textit{``#1''}}
\newcommand\norm[1]{\left\Vert#1\right\Vert}

\def\BibTeX{{\rm B\kern-.05em{\sc i\kern-.025em b}\kern-.08em
    T\kern-.1667em\lower.7ex\hbox{E}\kern-.125emX}}
\begin{document}

\title{Robust Precoding Using Channel Estimation Error Statistics for Cell-Free Massive MIMO in O-RAN\\}

\author{\IEEEauthorblockN{Mohammad Hossein Shokouhi and Vincent W.S. Wong}
\IEEEauthorblockA{Department of Electrical and Computer Engineering, The University of British Columbia, Vancouver, Canada \\
email: \{mhshokouhi, vincentw\}@ece.ubc.ca}
}

\maketitle
\begin{abstract}
The open radio access network (O-RAN) architecture enables hierarchical network control for cell-free massive multiple-input multiple-output (MIMO). Pilot assignment and precoding are tightly coupled components in cell-free massive MIMO systems, yet most existing works optimize them independently. Conventional precoding schemes often ignore pilot assignment decisions and treat estimated channels as perfect, which may lead to degraded performance under severe pilot contamination. 
In this paper, we propose a pilot assignment and robust precoding framework for cell-free massive MIMO in O-RAN. We formulate an optimization problem to maximize the aggregate ergodic rate under open radio unit (O-RU) transmit power and user minimum rate constraints. We use the block coordinate descent approach to alternate between pilot assignment and precoding optimization. 
For pilot assignment, we use multi-agent deep reinforcement learning (DRL), where a DRL agent is assigned to each user to update its pilot signal. 
For precoding, we derive a closed-form lower bound on the ergodic rate and propose an iterative optimization algorithm to determine the precoding vectors. Notably, the precoding update equation incorporates channel estimation error statistics to prioritize users with reliable channel estimates. 
Simulation results demonstrate that our proposed framework achieves higher aggregate throughput than three state-of-the-art baseline schemes. We also conduct ablation studies to show the individual contributions of pilot assignment and robust precoding components in our proposed framework.
\end{abstract}
\begin{IEEEkeywords}
Open radio access network (O-RAN), cell-free massive multiple-input multiple-output (MIMO), pilot assignment, precoding
\end{IEEEkeywords}
\section{Introduction}
The open radio access network (O-RAN) architecture enables disaggregated and hierarchical resource management in next-generation wireless networks \cite{survey:O-RAN1}. It introduces two RAN intelligent controllers (RICs) that operate at different timescales: the near-real-time (near-RT) RIC, which performs data-driven control via xApps, and the non-real-time (non-RT) RIC, which supports model training via rApps \cite{survey:O-RAN1}. In addition, O-RAN disaggregates the base station into the open central unit (O-CU), open distributed unit (O-DU), and open radio unit (O-RU), where real-time control logic can be implemented through dApps \cite{melodia:dApp_full}.

In conventional O-RAN, each user is served by a single O-RU. As a result, throughput degrades as the distance from the serving O-RU increases. Cell-free massive multiple-input multiple-output (MIMO) addresses this limitation by allowing multiple O-RUs to coherently serve each user \cite{Bjornson:CF-mMIMO}, which leads to more uniform throughput across the coverage area.

Reliable channel estimation is essential for signal processing in cell-free massive MIMO. Ideally, users transmit orthogonal pilot sequences to their serving O-RUs for channel estimation. However, if the number of users exceeds the number of available pilots, some users have to share pilots, which can lead to pilot contamination and channel estimation errors \cite{MIMO:pilot1}. Various pilot assignment algorithms have been proposed in the literature to reduce the channel estimation error.
In \cite{MIMO:pilot_approx_ratio}, semidefinite relaxation is used to solve the pilot assignment problem with the objective of minimizing pilot reuse.
In \cite{MIMO:pilot_quantum}, a hybrid quantum convolutional neural network (CNN) is proposed for pilot assignment to maximize the aggregate throughput.
In \cite{MIMO:pilot1}, a multi-agent deep reinforcement learning (DRL) algorithm updates the pilot signals in each near-RT loop to minimize the aggregate channel estimation error.
In \cite{MIMO:pilot_power}, random pilot assignment is considered. A multi-agent DRL algorithm is proposed to maximize the uplink throughput by optimizing the transmit powers of the pilot signals.
In \cite{MIMO:pilot_power2}, fractional programming is used for pilot power control and a multi-agent RL algorithm is employed for pilot assignment. 
% These works primarily focus on minimizing channel estimation errors rather than maximizing aggregate downlink throughput.

After channel estimation, the next step is precoding, where O-RUs steer the transmit signal beams towards the users they serve based on the channel estimates. Various precoding schemes have been proposed to maximize the aggregate user throughput.
In \cite{MIMO:MHSH1}, a multi-agent actor-critic DRL algorithm is proposed where an actor is assigned to each O-RU to locally determine its precoding matrices.
In \cite{MIMO:MHSH_TWC_2025}, the problem of maximizing the aggregate throughput subject to user minimum rate requirements is reformulated as a weighted minimum mean square error (WMMSE) problem and solved iteratively for cell-free massive MIMO systems.

Existing pilot assignment works \cite{MIMO:pilot1, MIMO:pilot_approx_ratio, MIMO:pilot_quantum, MIMO:pilot_power, MIMO:pilot_power2} typically rely on heuristic precoding schemes. Furthermore, existing precoding approaches \cite{MIMO:MHSH1, MIMO:MHSH_TWC_2025} focus solely on precoding under the assumption of perfect channel state information (CSI). The precoding algorithms in these works do not account for channel estimation errors, which may lead to performance degradation, especially in dense deployments with heavy pilot reuse. In this work, we propose a robust precoding scheme that incorporates pilot assignment decisions and channel estimation error statistics to prioritize users with reliable CSI. The contributions of this paper are summarized as follows:
\begin{itemize}
    \item We formulate a joint optimization problem for pilot assignment and precoding to maximize the aggregate ergodic throughput, subject to per-O-RU transmit power constraints and user minimum rate requirements. We adopt the BCD approach to decompose the problem into subproblems handled at different control timescales. In particular, the precoding vectors are optimized at the O-DU in real time, whereas the pilot assignment decisions are updated at the near-RT RIC using a multi-agent DRL-based framework.
    \item The ergodic-rate objective involves an expectation over channel uncertainty and does not admit a closed-form expression. To address this, we derive a closed-form lower bound on the ergodic rate and develop an iterative optimization algorithm to determine the precoding vectors. The proposed algorithm explicitly incorporates channel estimation error statistics to prioritize users with more reliable channel estimates. Such robust design reduces inter-user interference compared to conventional precoding approaches that treat channel estimates as perfect, particularly in dense deployments with heavy pilot reuse.
    \item We conduct extensive simulations to evaluate the proposed framework with different numbers of users and O-RUs. Results show that the proposed framework achieves higher aggregate throughput and lower channel estimation mean squared error (MSE) when compared with three state-of-the-art baselines. We also perform ablation studies to evaluate the impact of the pilot assignment and robust precoding components on the performance.
\end{itemize}

The rest of this paper is organized as follows. In Section \ref{sec:model}, we introduce the system model and formulate the optimization problem. In Section \ref{sec:algorithm}, we propose an iterative algorithm for robust precoding and a DRL algorithm for pilot assignment. Performance evaluation is presented in Section \ref{sec:sim}. Conclusion is given in Section \ref{sec:conclusion}.

{\it Notations}: In this paper, $\mathbb{C}$ denotes the set of complex numbers. We use boldface upper-case letters (e.g., $\mathbf{X}$) to denote matrices and boldface lower-case letters (e.g., $\mathbf{x}$) to denote vectors. $\mathbf{I}_N$ represents an $N \times N$ identity matrix. $(\cdot)^\textrm{H}$ denotes the conjugate transpose of a vector or a matrix.

\begin{figure}[t]
\center{\includegraphics[width=\linewidth]{Fig/ORAN/System_Model.pdf}}
\caption{The considered system model. rApps in the non-RT RIC perform training and update the DRL agent parameters. xApps in the near-RT RIC perform pilot assignment in each near-RT loop. dApps in the O-DU perform channel estimation and robust precoding.}
\label{fig:model}
\end{figure}

\section{System Model and Problem Formulation} \label{sec:model}
We consider downlink transmission in a cell-free massive MIMO system in O-RAN, as shown in Fig. \ref{fig:model}. Let $\mathcal{K}=\{1,2,\ldots,K\}$ denote the set of single-antenna users, where $K$ is the total number of users. We consider $L$ O-RUs denoted by set $\mathcal{L}=\{1,2,\ldots,L\}$. Each O-RU is equipped with $N_\textrm{t}$ transmit antennas. Let $\mathcal{K}_l$ and $K_l$ denote the set and the number of users served by O-RU $l$, respectively. Moreover, let $\mathcal{L}_k$ and $L_k$ denote the set and the number of O-RUs that serve user $k$, respectively. The system operates in time division duplexing (TDD) mode. Thus, uplink-downlink channel reciprocity holds.

We adopt a block fading model where the channel remains constant within each RT control loop and changes across RT loops \cite{Bjornson:CF-mMIMO}. Thus, each RT control loop corresponds to one coherence block.
Let $\mathbf{h}_{k,l}(t)=\sqrt{\beta_{k,l}}\,\mathbf{g}_{k,l} \in \mathbb{C}^{N_\textrm{t}}$ denote the channel vector between user $k \in \mathcal{K}$ and O-RU $l \in \mathcal{L}$, where $\beta_{k,l}$ denotes the large-scale fading coefficient from O-RU $l$ to user $k$ and $\mathbf{g}_{k,l}\in\mathbb{C}^{N_\textrm{t}}$ is the small-scale fading vector.

Let $\tau_\mathrm{c}$ denote the coherence block length, i.e., the total number of symbols over which the channel is constant. In each coherence block, the first $\tau_\mathrm{p}$ symbols are used for uplink pilot transmission and the remaining $\tau_\mathrm{d}=\tau_\mathrm{c}-\tau_\mathrm{p}$ symbols are used for downlink transmission. Thus, at most $\tau_\mathrm{p}$ mutually orthogonal pilots are available. Let $\boldsymbol{\Phi}=\{\boldsymbol{\phi}_1,\ldots,\boldsymbol{\phi}_{\tau_\mathrm{p}}\}$ denote the pilot codebook, where $\boldsymbol{\phi}_i\in\mathbb{C}^{\tau_\mathrm{p}}$ for $i=1,\ldots,\tau_\mathrm{p}$ is a unit-norm pilot sequence. We have $\boldsymbol{\phi}_i^\mathrm{H}\boldsymbol{\phi}_{i'}=1$ if $i=i'$ and $0$ otherwise. Let $\mathcal{T}_\mathrm{p}=\{1,\ldots,\tau_\mathrm{p}\}$ denote the set of pilot indices.

Let $a_{k,i}$ denote the pilot assignment variable, where $a_{k,i}=1$ if pilot $i\in\mathcal{T}_\mathrm{p}$ is assigned to user $k\in\mathcal{K}$, and $a_{k,i}=0$ otherwise. Thus, the pilot signal $\boldsymbol{\psi}_{k} \in \mathbb{C}^{\tau_\mathrm{p}}$ of user $k$ can be expressed as $\boldsymbol{\psi}_k=\sum_{i=1}^{\tau_\mathrm{p}} a_{k,i}\boldsymbol{\phi}_i$.

The received uplink signal $\mathbf{Y}_l^\textrm{UL} \in \mathbb{C}^{N_\textrm{t} \times \tau_\mathrm{p}}$ at O-RU $l$ is given by
\begin{equation}\label{eq:recv_uplink_signal}
    \mathbf{Y}_l^\textrm{UL}
    =
    \sum_{k\in\mathcal{K}} \sqrt{p_k}\,\mathbf{h}_{k,l}\boldsymbol{\psi}_k^\mathrm{H}
    + \mathbf{N}_l,
\end{equation}
where $p_k$ denotes the uplink pilot transmit power of user $k$. $\mathbf{N}_l \in \mathbb{C}^{N_\textrm{t} \times \tau_\mathrm{p}}$ is the additive white Gaussian noise matrix whose entries are independent and identically distributed (i.i.d.) complex Gaussian random variables with zero mean and variance $\left(\sigma^{\mathrm{UL}}\right)^2$, i.e., $\mathcal{CN}\left(0,\left(\sigma^{\mathrm{UL}}\right)^2\right)$.

To estimate $\mathbf{h}_{k,l}$, O-RU $l$ first correlates the received uplink signal $\mathbf{Y}_l^\textrm{UL}$ with the pilot signal $\boldsymbol{\psi}_{k}$ as $\mathbf{y}_{k,l}=\mathbf{Y}_l^\textrm{UL}\boldsymbol{\psi}_k$.
We next apply the linear minimum mean-square error (LMMSE) estimator. Let $\mathcal{P}_k\triangleq\{k':\boldsymbol{\psi}_{k'}=\boldsymbol{\psi}_k\}$ denote the set of users that share the same pilot as user $k$. Under the uncorrelated Rayleigh fading channel model, the LMMSE estimate of $\mathbf{h}_{k,l}$ based on the observation $\mathbf{y}_{k,l}$ is given by
\begin{equation}\label{eq:lmmse_final}
    \hat{\mathbf{h}}_{k,l}
    =
    \frac{\sqrt{p_k}\,\beta_{k,l}}
    {\sum\limits_{k'\in\mathcal{P}_k} p_{k'}\beta_{k',l} + \left(\sigma^{\mathrm{UL}}\right)^2}
    \mathbf{y}_{k,l}.
\end{equation}

The LMMSE estimate $\hat{\mathbf{h}}_{k,l}$ and the corresponding estimation error $\tilde{\mathbf{h}}_{k,l} \triangleq \mathbf{h}_{k,l} - \hat{\mathbf{h}}_{k,l}$ are independent and both follow complex Gaussian distributions. In particular, we have $\hat{\mathbf{h}}_{k,l} \sim \mathcal{CN}\left(\mathbf{0}, \alpha_{k,l}\mathbf{I}_{N_\mathrm{t}}\right)$ and $\tilde{\mathbf{h}}_{k,l} \sim \mathcal{CN}\left(\mathbf{0}, \left(\beta_{k,l} - \alpha_{k,l}\right)\mathbf{I}_{N_\mathrm{t}}\right)$, where $\alpha_{k,l}$ is given by
\begin{equation}
\alpha_{k,l}
=
\frac{p_k \beta_{k,l}^2}
{\sum\limits_{k'\in\mathcal{P}_k} p_{k'}\beta_{k',l} + \left(\sigma^{\mathrm{UL}}\right)^2}.
\end{equation}
Therefore, the estimation error covariance matrix is
\begin{equation}\label{eq:estimation_error_covariance}
\mathbf{R}_{k,l}
=
\mathbb{E}\left[\tilde{\mathbf{h}}_{k,l}\tilde{\mathbf{h}}_{k,l}^\mathrm{H}\right]
=
\left(\beta_{k,l} - \alpha_{k,l}\right)\mathbf{I}_{N_\mathrm{t}},
\end{equation}

The precoding vectors are designed based on the obtained channel estimates in each coherence block. The precoding vector at O-RU $l$ for user $k\in\mathcal{K}_l$ is denoted by $\mathbf{v}_{k,l} = \boldsymbol{\pi}_{k,l}\left(\hat{\mathbf{H}}\right) \in\mathbb{C}^{N_\textrm{t}}$, where $\hat{\mathbf{H}} = \left[\hat{\mathbf{h}}_{k,l},\;\forall k \in \mathcal{K}, l \in \mathcal{L}\right] \in \mathbb{C}^{LN_\textrm{t} \times K}$ and $\boldsymbol{\pi}_{k,l}(\cdot)$ is the precoding policy.

The downlink signal transmitted by O-RU $l\in\mathcal{L}$ is $\mathbf{x}_l=\sum_{k\in\mathcal{K}_l}\mathbf{v}_{k,l}s_k$,
% \begin{equation}\label{eq:x_l}
%     \mathbf{x}_l=\sum_{k\in\mathcal{K}_l}\mathbf{v}_{k,l}s_k,
% \end{equation}
where $s_k\in\mathbb{C}$ is the data symbol for user $k$, with $\mathbb{E}[|s_k|^2]=1$ and $\mathbb{E}[s_k s_i^*]=0$ for $i\neq k$. The signal received by user $k$ can be written as $    y_k^\textrm{DL} = \sum_{l\in\mathcal{L}} \mathbf{h}_{k,l}^\mathrm{H}\mathbf{x}_l + n_k$,
% \begin{equation}\label{eq:y_k}
%     y_k^\textrm{DL}
%     =
%     \sum_{l\in\mathcal{L}} \mathbf{h}_{k,l}^\mathrm{H}\mathbf{x}_l + n_k,
% \end{equation}
where $n_k\sim\mathcal{CN}\left(0,\left(\sigma^{\mathrm{DL}}\right)^2\right)$ is additive white Gaussian noise.
The ergodic downlink rate of user $k$ is given by
\begin{equation}\label{eq:rate_uatf_start}
    R_k^{\textrm{DL}}
    =
    \frac{\tau_{\textrm{d}}}{\tau_{\textrm{c}}} \mathbb{E}\left[\log_2\left(1+\gamma_k^\textrm{DL}\right)\right],
\end{equation}
where $\gamma_k^\textrm{DL}$ is the effective downlink signal-to-interference-plus-noise ratio (SINR) of user $k$, given by
\begin{equation}\label{eq:SINR}
    \gamma_k^\textrm{DL}
    =
    \frac{\left|\xi_{k,k}\right|^2}
    {\sum_{i\in\mathcal{K} \setminus\{k\}}\left|\xi_{k,i}\right|^2
    +\left(\sigma^{\mathrm{DL}}\right)^2 },
\end{equation}
In \eqref{eq:SINR}, $\xi_{k,i} = \sum_{l\in \mathcal{L}_i}\mathbf{h}_{k,l}^\mathrm{H}\mathbf{v}_{i,l} \in \mathbb{C}$ denotes the effective channel for user $k$ if $i=k$, and the effective interference from user $i$ to user $k$ otherwise.

We aim to maximize the aggregate downlink throughput of users. The optimization problem is
\begin{subequations} \label{eq:optimization}
\begin{alignat}{2}
&\maximize\limits_{\substack{
\mathbf{v}_{k,l},\,a_{k,i}\\
k\in\mathcal{K},l\in\mathcal{L},i\in\mathcal{T}_{\mathrm{p}}
}}
& &\quad \sum_{k\in\mathcal{K}} R_k^{\mathrm{DL}} \label{eq:obj}\\
&\textrm{subject to}
& &\quad R_k^{\mathrm{DL}} \ge R_k^{\min},\; k\in\mathcal{K} \label{eq:cons_rate_min}\\
& & &\quad \sum_{k\in\mathcal{K}_l}\|\mathbf{v}_{k,l}\|_2^2 \le P_\textrm{DL}^{\textrm{max}},\; l\in\mathcal{L} \label{eq:cons_dl_power}\\
% & & &\quad \|\mathbf{v}_{k,l}\|_2^2 \le c_{k,l}P_\textrm{DL}^{\textrm{max}},\; k\in\mathcal{K},\,l\in\mathcal{L} \label{eq:cons_link_cv}\\
& & &\quad \sum_{i\in\mathcal{T}_{\mathrm{p}}} a_{k,i} = 1,\; k\in\mathcal{K} \label{eq:cons_one_pilot}\\
& & &\quad a_{k,i}\in\{0,1\},\; k\in\mathcal{K},\; i\in\mathcal{T}_{\mathrm{p}}. \label{eq:cons_binary_a}
\end{alignat}
\end{subequations}

Constraint \eqref{eq:cons_rate_min} guarantees the minimum rate requirement $R_k^{\min}$ for each user $k\in\mathcal{K}$.
Constraint \eqref{eq:cons_dl_power} limits the total downlink transmit power at each O-RU to $P_\textrm{DL}^{\textrm{max}}$.
Constraint \eqref{eq:cons_one_pilot} enforces that each user is assigned exactly one pilot sequence.
Problem \eqref{eq:optimization} is a mixed-integer optimization problem. It is NP-hard and difficult to solve jointly in real time.

\section{DRL-Based Pilot Assignment and Robust Precoding for Cell-free Massive MIMO in O-RAN} \label{sec:algorithm}
We decompose problem \eqref{eq:optimization} using BCD. The precoding vectors are optimized at the O-DUs based on instantaneous channel estimates, while the pilot assignment is updated at the near-RT RIC. We first present the robust precoding algorithm, followed by the DRL-based pilot assignment algorithm.
\subsection{Precoding Subproblem}
For fixed pilot assignment, the precoding subproblem remains intractable because the ergodic rate in \eqref{eq:rate_uatf_start} does not admit a closed-form expression. Therefore, we propose the following tractable lower bound.
\begin{lemma}\label{lemma:ergodic_lower_bound}
A lower bound on the ergodic downlink rate of user $k$ is given by
\begin{equation}\label{eq:ergodic_bound}
    R_k^{\textrm{DL}} \geq \mathbb{E}_{\hat{\mathbf{H}}} \left[\hat{R}_k^{\textrm{DL}}\right],
\end{equation}
where
\begin{equation}\label{eq:conditional_UatF}
    \hat{R}_k^{\textrm{DL}} = \frac{\tau_{\textrm{d}}}{\tau_{\textrm{c}}} \log_2\left(1+\hat{\gamma}_k^{\textrm{DL}}\right),
\end{equation}
and $\hat{\gamma}_k^{\textrm{DL}}$ is given by
\begin{align}
\hat{\gamma}_k^{\textrm{DL}} =
\frac{\left|\hat{\xi}_{k,k}\right|^2}
{\sum\limits_{i\in\mathcal{K} \setminus \{k\}} \left|\hat{\xi}_{k,i}\right|^2
+ \sum\limits_{i \in \mathcal{K}} \sum\limits_{l\in \mathcal{L}_i}
\mathbf{v}_{i,l}^\mathrm{H} \mathbf{R}_{k,l} \mathbf{v}_{i,l}
+ \left(\sigma^{\mathrm{DL}}\right)^2},
\end{align}
with $\hat{\xi}_{k,i} = \sum_{l\in \mathcal{L}_i} \hat{\mathbf{h}}_{k,l}^\mathrm{H}\mathbf{v}_{i,l}$.
\end{lemma}

\begin{IEEEproof}
    The proof is provided in Appendix \ref{appendix:ergodic_lower_bound}.
\end{IEEEproof}

The lower bound in Lemma \ref{lemma:ergodic_lower_bound} admits a closed-form expression, which enables optimization with respect to the precoding variables. We replace the ergodic rate in the precoding subproblem with its lower bound as
\begin{subequations} \label{eq:precoding_subproblem}
\begin{alignat}{2}
&\maximize\limits_{\substack{
\mathbf{v}_{k,l},\\
k \in \mathcal{K}_l, l \in \mathcal{L}
}}
& &\quad \sum_{k\in\mathcal{K}} \mathbb{E}_{\hat{\mathbf{H}}} \left[\hat{R}_k^{\textrm{DL}}\right] \label{eq:precoding_obj}\\
&\textrm{subject to}
& &\quad \mathbb{E}_{\hat{\mathbf{H}}} \left[\hat{R}_k^{\textrm{DL}}\right] \ge R_k^{\min},\; k\in\mathcal{K} \label{eq:precoding_cons_rate_min}\\
& & &\quad \textrm{constraint \eqref{eq:cons_dl_power}}. \notag
\end{alignat}
\end{subequations}

To incorporate the minimum data rate constraint \eqref{eq:precoding_cons_rate_min} into the objective function, we introduce the Lagrange multipliers $\boldsymbol{\mu}=[\mu_1,\ldots,\mu_K]^\top$. For fixed $\boldsymbol{\mu}$, the resulting weighted precoding subproblem is
\begin{align}\label{eq:precoding_subproblem_lagrangian}
&\maximize\limits_{\substack{
\mathbf{v}_{k,l},\\
k \in \mathcal{K}_l,\; l \in \mathcal{L}
}}
\quad \sum_{k\in\mathcal{K}} \eta_k \mathbb{E}_{\hat{\mathbf{H}}} \left[\hat{R}_k^{\textrm{DL}}\right]\\
&\textrm{subject to}
\quad \textrm{constraint \eqref{eq:cons_dl_power}}, \notag
\end{align}
where $\eta_k = 1+\mu_k, \; k\in\mathcal{K}$.
Since the precoding vectors are designed in each coherence block using the instantaneous channel estimates, problem \eqref{eq:precoding_subproblem_lagrangian} can be decoupled across coherence blocks as
\begin{align}\label{eq:precoding_subproblem_instantaneous}
&\maximize\limits_{\substack{
\mathbf{v}_{k,l},\\
k \in \mathcal{K}_l,\; l \in \mathcal{L}
}}
\quad \sum_{k\in\mathcal{K}} \eta_k \hat{R}_k^{\mathrm{DL}}\left(\hat{\mathbf{H}}\right)\\
&\textrm{subject to}
\quad \textrm{constraint \eqref{eq:cons_dl_power}}. \notag
\end{align}

Problem \eqref{eq:precoding_subproblem_instantaneous} is still nonconvex due to the nonconvexity of $\hat{R}_k^{\mathrm{DL}}\left(\hat{\mathbf{H}}\right)$. Using the receive filter $u_k \in \mathbb{C}$, the data symbol estimated by user $k$ is $\hat{s}_k = u_k^* y_k^\textrm{DL}$. Consequently, the conditional MSE of the symbol estimated by user $k$ is
\begin{align}
e_k \!=\!{}&\mathbb{E}\!\left[
\left| s_k - \hat{s}_k \right|^2 \mid \hat{\mathbf{H}}
\right]
\!=\!\left|1 - u_k^*\hat{\xi}_{k,k}\right|^2 \!+\! \sum\limits_{i\in\mathcal{K}\setminus\{k\}} \left|u_k\right|^2\left|\hat{\xi}_{k,i}\right|^2 \nonumber \\
&+ \left|u_k\right|^2 \sum\limits_{i \in \mathcal{K}} \sum\limits_{l\in \mathcal{L}_i}
\mathbf{v}_{i,l}^\mathrm{H} \mathbf{R}_{k,l} \mathbf{v}_{i,l}
+ \left|u_k\right|^2 \left(\sigma^{\mathrm{DL}}\right)^2.
\end{align}
By introducing the auxiliary weight $w_k \geq 0$ for each user $k$, problem \eqref{eq:precoding_subproblem_instantaneous} can be reformulated as an equivalent WMMSE problem, expressed as
\begin{align}
&\minimize_{\substack{
w_k, u_k, \mathbf{v}_{k,l},\\
k \in \mathcal{K},\; l \in \mathcal{L}_k
}} \quad\sum_{k \in \mathcal{K}} \eta_{k}\left(w_k e_k - \ln  (w_k)\right) \label{eq:WMMSE_problem}\\
&\textrm{subject to constraint \eqref{eq:cons_dl_power}}. \notag
\end{align} 

Problem \eqref{eq:WMMSE_problem} is convex with respect to each of the individual optimization variables $u_k, w_k, \textrm{ and } \mathbf{v}_{k,l}$.
By fixing $\mathbf{v}_{k,l}$, the optimal $u_k^\mathrm{opt}$ can be obtained as \cite{MIMO1}
\begin{equation}\label{eq:u_WMMSE}
    u_k^\mathrm{opt} = \frac{\hat{\xi}_{k,k}}{\sum_{i \in \mathcal{K}} \left|\hat{\xi}_{k,i}\right|^2 + 
    \sum\limits_{i \in \mathcal{K}} \sum\limits_{l\in \mathcal{L}_i}
    \mathbf{v}_{i,l}^\mathrm{H} \mathbf{R}_{k,l} \mathbf{v}_{i,l} + 
    \left(\sigma^{\mathrm{DL}}\right)^2
    }.
\end{equation}
Moreover, by inserting the optimal $u_k^\mathrm{opt}$ and fixing $\mathbf{v}_{k,l}$, the optimal $w_k^\mathrm{opt}$ can be determined as
\begin{equation}\label{eq:w_WMMSE}
    w_k^\mathrm{opt} = \frac{1}{e_k^\mathrm{opt}},
\end{equation}
where
\begin{equation}
    e_k^\mathrm{opt} = 1 - \frac{\left|\hat{\xi}_{k,k}\right|^2}{\sum_{i \in \mathcal{K}} \left|\hat{\xi}_{k,i}\right|^2 + 
    \sum\limits_{i \in \mathcal{K}} \sum\limits_{l\in \mathcal{L}_i}
    \mathbf{v}_{i,l}^\mathrm{H} \mathbf{R}_{k,l} \mathbf{v}_{i,l} + 
    \left(\sigma^{\mathrm{DL}}\right)^2
    }.
\end{equation}
Finally, by holding $u_k$ and $w_k$ fixed, problem \eqref{eq:WMMSE_problem} can be solved with respect to $\mathbf{v}_{k,l}$. Note that the resulting precoding subproblem cannot be decoupled across O-RUs. This is because, in cell-free massive MIMO with coherent joint transmission, multiple O-RUs collaboratively serve each user, which couples their precoding vectors \cite{MIMO:MHSH_TWC_2025}.
To address this, we employ the BCD method \cite{MIMO:MHSH_TWC_2025}, where the precoding vectors $\mathbf{v}_{k,l}$ for each O-RU $l$ are updated iteratively while keeping the precoding vectors of the other O-RUs fixed. Using this approach, the update equation for $\mathbf{v}_{k,l}$ can be derived as
\begin{align}\label{eq:V_WMMSE}
    \mathbf{v}_{k,l}^\mathrm{opt} = &\left(
    \sum_{i \in \mathcal{K}} \eta_i w_i |u_i|^2 \left(\hat{\mathbf{h}}_{i,l} \hat{\mathbf{h}}_{i,l}^\mathrm{H} + \mathbf{R}_{i,l}\right)
    + \lambda_l \mathbf{I}_{N_\mathrm{t}}
    \right)^{-1} \nonumber \\
    & \left(
    \eta_k w_k u_k \hat{\mathbf{h}}_{k,l} - \sum_{i \in \mathcal{K}} \eta_i w_i |u_i|^2 \hat{\mathbf{h}}_{i,l} z_{i,k,l}^* 
    \right),
\end{align}
where $z_{i,k,l} \triangleq \sum_{j \in \mathcal{L}_k \setminus\{l\}} \hat{\mathbf{h}}_{i,j}^\mathrm{H} \mathbf{v}_{k,j}$ and $\lambda_l \geq 0$ is a Lagrange multiplier that must satisfy the complementary slackness condition of constraint \eqref{eq:cons_dl_power}.


Finally, $\mu_k$ is updated using projected gradient ascent as
\begin{equation}\label{eq:mu_update}
    \mu_k \xleftarrow{} \left[\mu_k + \delta_k \left( R_{k}^{\textrm{min}} - \mathbb{E}_{\hat{\mathbf{H}}} \left[\hat{R}_k^{\textrm{DL}}\right] \right)\right]^+, \quad k \in \mathcal{K},
\end{equation}
where $\delta_k$ is the step size for user $k$. 

The proposed precoding algorithm explicitly incorporates the estimation error covariance matrix $\mathbf{R}_{k,l}$. Users with unreliable CSI are penalized through the additional error-dependent terms in \eqref{eq:u_WMMSE}, \eqref{eq:w_WMMSE}, and \eqref{eq:V_WMMSE}. Unlike conventional precoding schemes that treat channel estimates as perfect, the proposed framework adapts the power allocation based on estimation error statistics and prioritizes users with more reliable CSI, which reduces interference caused by inaccurate precoding.

%% Section III: DRL-Based Pilot Assignment
%% To be compiled as part of the main document.

\subsection{Pilot Assignment Subproblem}\label{sec:pilot_assignment}
The pilot assignment subproblem is combinatorial with search space $\tau_\mathrm{p}^K$. To solve it efficiently, we model it as a Markov game, where a DRL agent is assigned to each user $k \in \mathcal{K}$ to determine its pilot signal. 
Each near-RT loop is treated as one decision step of the Markov game \cite{MIMO:pilot1}. Each non-RT loop corresponds to a training episode.
To avoid non-stationarity caused by simultaneous updates, we adopt a sequential strategy where only the pilot signal of a single user is updated in each near-RT loop.

In each near-RT loop, a target user $k^\star$ is selected by calculating a priority score for each user $k \in \mathcal{K}$ as
\begin{equation}\label{eq:priority_score}
    \rho_k = \omega_\mathrm{R} \left[R_k^{\min} - R_k^{\mathrm{DL}}\right]_+ + \omega_\mathrm{C} \, C_k + \omega_\mathrm{U} \, U_k,
\end{equation}
where $\omega_\mathrm{R}, \omega_\mathrm{C}, \omega_\mathrm{U} \geq 0$ are weighting coefficients. The term $\left[R_k^{\min} - R_k^{\mathrm{DL}}\right]_+$ represents the current minimum rate requirement violation of user $k$. The pilot conflict term $C_k = \sum_{i \in \mathcal{P}_k\setminus\{k\}} \chi_{k,i}$ indicates the severity of pilot contamination, where $\chi_{k,i} = \sum_{l \in \mathcal{L}_k \cap \mathcal{L}_i} \beta_{k,l} \, \beta_{i,l}$.
% \begin{equation}\label{eq:conflict_weight}
%     \chi_{k,i} = \sum_{l \in \mathcal{L}_k \cap \mathcal{L}_i} \beta_{k,l} \, \beta_{i,l}.
% \end{equation}
The CSI uncertainty term is defined as $U_k = \sum_{l \in \mathcal{L}_k} \tr \left(\mathbf{R}_{k,l}\right)$. The user with the highest priority score is selected as $k^\star = \arg\max_{k \in \mathcal{K}} \, \rho_k$.

Once $k^\star$ is selected, the DRL agent assigned to user $k^\star$ updates the pilot signal of that user.
Each agent $k$ has an observation space $\mathcal{O}_k$, an action space $\mathcal{A}_k$, a policy $\pi_k$, and a reward function $r_k$.
The observation $\mathbf{o}_{k^\star}$ of agent $k^\star$ includes the current pilot index of user $k^\star$, its ergodic rate $R_{k^\star}^{\textrm{DL}}$, QoS deficit $\left[R_{k^\star}^{\min} - R_{k^\star}^{\textrm{DL}}\right]_+$, and CSI uncertainty $U_{k^\star}$. 
In addition, it includes the pilot conflict scores $c_{k^\star, t} \triangleq \sum_{i \in \mathcal{K} \setminus \{k^\star\} : \boldsymbol{\psi}_i = \boldsymbol{\phi}_t} \chi_{k^\star, i},\; t \in \mathcal{T}_\mathrm{p}$ as well as the pilot occupancy vector $\mathbf{n} = [n_1, \ldots, n_{\tau_\mathrm{p}}]^\top$, where $n_t = \left|\{i \in \mathcal{K} : \boldsymbol{\psi}_i = \boldsymbol{\phi}_t\}\right|$. 
To provide global context for agent $k^\star$, we define $\mathcal{M}_{k^\star}$, which is the set of $M$ users with highest overlap weights $\chi_{k^\star,m}$. Consequently, the observation $\mathbf{o}_{k^\star}$ also includes the current pilot index $\boldsymbol{\psi}_m$, overlap weight $\chi_{k^\star,m}$, and QoS deficit $\left[R_{m}^{\min} - R_{m}^{\textrm{DL}}\right]_+$ for $m \in \mathcal{M}_{k^\star}$.
The action $a_{k^\star} \in \mathcal{T}_\mathrm{p} = \{1, \ldots, \tau_\mathrm{p}\}$ is the new pilot index assigned to user $k^\star$. 

The reward function must reflect system-wide performance to discourage selfish policies. It is defined as the change in the objective function $J=\sum_{k\in\mathcal{K}} \eta_k R_k^{\mathrm{DL}}$, expressed as
\begin{equation}\label{eq:reward}
    r_{k^\star} = J\left(\mathbf{s}'\right) - J(\mathbf{s}),
\end{equation}
where $\mathbf{s}$ and $\mathbf{s}'$ denote the environment state before and after applying action $a_{k^\star}$, respectively.

We solve the above Markov game using the Dueling Deep Q-Network (DQN) algorithm. The Dueling DQN approximates the action-value function $Q(\mathbf{o}_{k^\star}, a_{k^\star})$ by decomposing it into a state-value function $V\!\left(\mathbf{o}_{k^\star}; \boldsymbol{\theta}_V\right)$ and an advantage function $A\!\left(\mathbf{o}_{k^\star}, a_{k^\star}; \boldsymbol{\theta}_A\right)$,
% \begin{multline}\label{eq:Q_dueling}
%     Q\!\left(\mathbf{o}_{k^\star}, a_{k^\star}; \boldsymbol{\theta}\right)
%     =
%     V\!\left(\mathbf{o}_{k^\star}; \boldsymbol{\theta}_V\right)
%     +
%     \Bigg(
%     A\!\left(\mathbf{o}_{k^\star}, a_{k^\star}; \boldsymbol{\theta}_A\right)\\
%     -
%     \frac{1}{\tau_\mathrm{p}}
%     \sum_{a'=1}^{\tau_\mathrm{p}} A\!\left(\mathbf{o}_{k^\star}, a'; \boldsymbol{\theta}_A\right)
%     \Bigg),
% \end{multline}
where $\boldsymbol{\theta} = \{\boldsymbol{\theta}_V, \boldsymbol{\theta}_A\}$ denotes the learnable parameters. The value function $V\!\left(\mathbf{o}_{k^\star}; \boldsymbol{\theta}_V\right)$ estimates the expected long-term return from observation $\mathbf{o}_{k^\star}$ regardless of the selected action, while the advantage function $A\!\left(\mathbf{o}_{k^\star}, a_{k^\star}; \boldsymbol{\theta}_A\right)$ estimates the relative benefit of selecting action $a_{k^\star}$ over the average.
The Q-network consists of a shared feature extraction module followed by two output heads for estimating $V\!\left(\mathbf{o}_{k^\star}; \boldsymbol{\theta}_V\right)$ and $A\!\left(\mathbf{o}_{k^\star}, a_{k^\star}; \boldsymbol{\theta}_A\right)$. To improve sample efficiency and convergence speed, all agents share the same network parameters $\boldsymbol{\theta}$.

Given the observation $\mathbf{o}_{k^\star}$, the agent selects the pilot index according to $a_{k^\star} = \arg\max_{a \in \mathcal{T}_\mathrm{p}} Q\!\left(\mathbf{o}_{k^\star}, a; \boldsymbol{\theta}\right)$.
To stabilize training, we employ experience replay. Transitions $\left(\mathbf{o}_{k^\star}, a_{k^\star}, r, \mathbf{o}'_{k^\star}\right)$ are stored in a replay buffer $\mathcal{B}$, where $\mathbf{o}'_{k^\star}$ is the observation of agent ${k^\star}$ after taking action $a_{k^\star}$. Parameters $\boldsymbol{\theta}$ are optimized by minimizing the temporal-difference (TD) loss, defined as
\begin{equation}\label{eq:TD_loss}
    \mathcal{L}(\boldsymbol{\theta})
    =
    \mathbb{E}_{\left(\mathbf{o}_{k^\star}, a, r, \mathbf{o}'_{k^\star}\right) \sim \mathcal{B}}
    \!\left[
    \left(
    y - Q\!\left(\mathbf{o}_{k^\star}, a_{k^\star}; \boldsymbol{\theta}\right)
    \right)^2
    \right],
\end{equation}
where the TD target $y = r_{k^\star} + \zeta \max_{a'} Q\!\left(\mathbf{o}'_{k^\star}, a'; \bar{\boldsymbol{\theta}}\right)$, $\zeta$ is the discount factor, and $\bar{\boldsymbol{\theta}}$ are the parameters of the target network. An $\varepsilon$-greedy policy is used during training to balance exploration and exploitation, where $\varepsilon$ is gradually annealed over the course of training. 
Training is performed offline in the non-RT RIC via rApps. After convergence, the learned parameters are transferred to the near-RT RIC and deployed within the xApps, as shown in Fig \ref{fig:model}.

A key advantage of the proposed framework is its adaptability to dynamic network conditions. Since the priority score $\rho_k$ in \eqref{eq:priority_score} is updated at each near-RT loop, users experiencing performance degradation are promptly selected for pilot reassignment. Furthermore, changes in large-scale fading coefficients are reflected in the observation $\mathbf{o}_{k^\star}$ through conflict scores and rate statistics, so the DRL agent updates its action $a_{k^\star}$ accordingly. As a result, the learned policy dynamically adapts to user mobility and time-varying channel conditions without requiring retraining.

\section{Performance Evaluation} \label{sec:sim}
In this section, we evaluate the performance of our proposed pilot assignment and robust precoding framework and compare it with three baseline schemes.
The considered cell-free O-RAN consists of $K=16$ users and $L=36$ O-RUs deployed in a $500 \textrm{ m}^2$ area. 
Each user is served by the $L^\textrm{max}=8$ O-RUs with the largest $\beta_{k,l}$.
The maximum transmit power of each O-RU $P^{\textrm{max}}$ is set to $30$ dBm.
The noise power is set to $\sigma^2= -114$ dBm.
Each O-RU is equipped with $N_{\textrm{t}}=4$ transmit antennas.
We use the uncorrelated Rayleigh fading channel model.
The large-scale fading coefficients $\beta_{k,l}$ follow the 3GPP urban microcell non-line-of-sight (UMi-NLOS) pathloss model \cite{3GPP_pathloss} with a carrier frequency of $f_\textrm{c}=2 \textrm{ GHz}$.
We set the duration of each RT loop $T=1$ ms. A non-RT loop occurs once every $N_\textrm{nRT}=10$ near-RT loops, and each near-RT loop comprises $N_\textrm{RT}=10$ RT loops.

\section{Conclusion} \label{sec:conclusion}
In this paper, we proposed a pilot assignment and robust precoding framework for cell-free massive MIMO in O-RAN. We formulated an optimization problem to maximize the aggregate ergodic rate under open radio unit (O-RU) transmit power and user minimum rate constraints. We used the BCD approach to alternate between pilot assignment and precoding. 
For pilot assignment, we proposed a multi-agent DRL algorithm that updates the pilot signal of a selected user in each near-RT loop. 
For precoding, we derived a closed-form lower bound on the ergodic rate and proposed an iterative optimization algorithm.
The precoding update equation incorporates channel estimation error covariance matrices to penalize users with unreliable CSI, which leads to reduced interference. 
Simulation results showed that our proposed framework achieves higher aggregate throughput than state-of-the-art baseline schemes. Ablation studies confirmed the effectiveness of the robust precoding component.
As future work, we aim to consider additional optimization variables, such as pilot power allocation.

\appendices
\section{Proof of Lemma \ref{lemma:ergodic_lower_bound}}
\label{appendix:ergodic_lower_bound}
\begin{IEEEproof}
Conditioned on $\hat{\mathbf{H}}$, the received signal of user $k$ can be written as
\begin{equation}
    y_k^{\mathrm{DL}}
    =
    \mathbb{E}\!\left[\xi_{k,k}\mid\hat{\mathbf{H}}\right]s_k
    + \tilde{n}_k,
\end{equation}
where
\begin{equation}
    \tilde{n}_k =
    \left(\xi_{k,k}-\mathbb{E}\!\left[\xi_{k,k}\mid\hat{\mathbf{H}}\right]\right)s_k
    + \sum_{i\in\mathcal{K}\setminus\{k\}}\xi_{k,i}s_i+n_k .
\end{equation}
Since $\tilde{n}_k$ is uncorrelated with $s_k$, treating $\tilde{n}_k$ as worst-case Gaussian noise gives the conditional lower bound
\begin{equation}
\frac{\tau_{\mathrm{d}}}{\tau_{\mathrm{c}}}
\mathbb{E}\!\left[
\log_2(1+\gamma_k^{\mathrm{DL}})
\mid\hat{\mathbf{H}}
\right]
\ge
\frac{\tau_{\mathrm{d}}}{\tau_{\mathrm{c}}}
\log_2\!\left(
1+
\frac{
\left|\mathbb{E}\!\left[\xi_{k,k}\mid\hat{\mathbf{H}}\right]\right|^2
}{
\mathbb{E}\!\left[|\tilde{n}_k|^2\mid\hat{\mathbf{H}}\right]
}
\right).
\end{equation}
The effective noise power can be written as
\begin{equation}
\begin{aligned}
\mathbb{E}\!\left[|\tilde{n}_k|^2\mid\hat{\mathbf{H}}\right]
=&
\sum_{i\in\mathcal{K}}
\mathbb{E}\!\left[|\xi_{k,i}|^2\mid\hat{\mathbf{H}}\right]
-
\left|\mathbb{E}\!\left[\xi_{k,k}\mid\hat{\mathbf{H}}\right]\right|^2 \\
&+
\left(\sigma^{\mathrm{DL}}\right)^2 .
\end{aligned}
\end{equation}
Using $\mathbf{h}_{k,l}=\hat{\mathbf{h}}_{k,l}+\tilde{\mathbf{h}}_{k,l}$,
$\mathbb{E}\left[\tilde{\mathbf{h}}_{k,l}\mid\hat{\mathbf{H}}\right]=\mathbf{0}$, and
$\mathbb{E}\left[\tilde{\mathbf{h}}_{k,l}\tilde{\mathbf{h}}_{k,l}^{\mathrm{H}}\mid\hat{\mathbf{H}}\right]=\mathbf{R}_{k,l}$, we obtain
\begin{equation}
    \mathbb{E}\!\left[\xi_{k,i}\mid\hat{\mathbf{H}}\right]
    =
    \hat{\xi}_{k,i},
\end{equation}
\begin{equation}
    \mathbb{E}\!\left[|\xi_{k,i}|^2\mid\hat{\mathbf{H}}\right]
    =
    |\hat{\xi}_{k,i}|^2
    +
    \sum_{l\in\mathcal{L}_i}
    \mathbf{v}_{i,l}^{\mathrm{H}}\mathbf{R}_{k,l}\mathbf{v}_{i,l}.
\end{equation}
Substituting these expressions into the conditional lower bound yields
$\hat{R}_k^{\mathrm{DL}}$ in \eqref{eq:conditional_UatF}. Finally, taking the expectation over $\hat{\mathbf{H}}$ and applying the law of total expectation gives
\begin{equation}
    R_k^{\mathrm{DL}}
    \ge
    \mathbb{E}_{\hat{\mathbf{H}}}
    \left[\hat{R}_k^{\mathrm{DL}}\right],
\end{equation}
which completes the proof.
\end{IEEEproof}
\bibliographystyle{IEEEtran}
\bibliography{Reference}
\end{document}